
\emph{Constraint results on the Presentation-A selection of Equation~\eqref{eq:V}.} This paragraph records constraint results bearing on the selection of the Higgs potential by Presentation~A, the missing input of Section~\ref{sec:ch3-relate}; no derivation of that selection from Equation~\eqref{eq:Smatrix} is obtained, and Assumption~\ref{ass:dictionary} remains the assumed matching condition throughout. Write \(X=\Phi^\dagger\Phi\). (i) Forced form and exact slice: under the stated bifundamental field content and the quartic degree-\((2,2)\) truncation, invariance under independent left and right unitary transformations of \(\Phi\) leaves exactly two real quartic invariants, \(\mathrm{Tr}\,X^2\) and \((\mathrm{Tr}\,X)^2\), with \(\mathrm{Tr}\,X^2=\sum_{i=1}^4\sigma_i(\Phi)^4\) in the singular values \(\sigma_i(\Phi)\) (hence non-negative), so the general quartic matter potential is the two-parameter family \(\mathcal V=c_3\,\mathrm{Tr}\,X^2+c_2(\mathrm{Tr}\,X)^2+c_1\,\mathrm{Tr}\,X+c_0\); Equation~\eqref{eq:V} is exactly the slice \((c_3,c_2,c_1,c_0)=(\lambda/4,\,0,\,-\lambda v^2/2,\,\lambda v^4)\) at internal dimension four, its constant \(c_0=\lambda v^4\) being the sum of the four squared terms; the selection question reduces to the ratio \(c_2/c_3\), which Equation~\eqref{eq:V} sets to zero, and to the signs. Every point of Equation~\eqref{eq:vacuum-manifold} is stationary for this potential, whose gradient \((\lambda/2)\,\Phi(X-v^2\mathbb I)\) is the potential term of Equation~\eqref{eq:Phi-eom}. (ii) Family theorem and slice threshold: throughout the strictly bounded-below region \(c_3>0\), \(c_3+4c_2>0\) - where the Hessian in \(x_i=\sigma_i(\Phi)^2\) has eigenvalues \(2c_3\) (three times) and \(2c_3+8c_2\), so the region is also strict convexity - the family has a unique vacuum \(X=s\,\mathbb I\) with \(s=-c_1/(2(c_3+4c_2))\); the infimum over \(\{\det\Phi=0\}\) is attained at singular values \((0,w,w,w)\) with \(w^2=-c_1/(2(c_3+3c_2))\); and the defect gap \(\Delta=(c_1^2/4)\,(c_3/12)/[(c_3/4+c_2)(c_3/3+c_2)]\) is strictly positive through the whole region. On the slice of Equation~\eqref{eq:V} one has \(s=v^2\), the vacuum is the manifold of Equation~\eqref{eq:vacuum-manifold}, and \(\Delta=\lambda v^4/4\): the defect threshold recorded above is re-derived as a slice of a family theorem, its positivity is region-wide while its value is slice-dependent, and off the slice the face minimizer moves, \(w\ne s\). (iii) Pinned truncation order: for the constant-\(\Theta\) Moyal product of Section~\ref{sec:ch3-relate} the two-point phase vanishes exactly and the four-point phase is cyclic on the momentum-conservation shell, while the quartic vertex carries the three channel factors \(\cos(\Theta p/2)\), \(\cos(\Theta q/2)\), \(\cos(\Theta(p+q)/2)\), pairwise distinct and coinciding flatly at \(\Theta=0\) with \(1-\cos(\Theta p/2)=\Theta^2p^2/8+O(\Theta^4)\); no choice of \(c_0,c_1,c_2,c_3\) absorbs these phases, so a channel-blind local \(\mathrm{Tr}\)-polynomial vertex such as Equation~\eqref{eq:V} can match the Presentation-A quartic vertex only through order \(\Theta^2\) - a truncation output, not an exact finite-\(\Theta\) output, of the assumed dictionary. (iv) Quartic-order boundary: beyond quartic order the gap property fails, since \(\varepsilon(\mathrm{Tr}\,X)^4-\mathrm{Tr}\,X^2\) with \(\varepsilon>0\) is bounded below with its global minimum on \(\{\det\Phi=0\}\) at rank one, of value \(-1/(4\varepsilon)\); the gap statements of (ii) are theorems of the quartic class, not of bounded matrix potentials. Scope: (i)-(iv) are identities and exact reductions, a family theorem conditional on the stated field content and truncation, and an explicit counterexample; the field content, the quartic truncation, constant \(\Theta\), and Assumption~\ref{ass:dictionary} are assumed and not derived, the selection of Equation~\eqref{eq:V} is constrained but not derived-closed, and nothing in this paragraph closes Problem~D, Problem~D1, or Problem~A.
